\scanpage{29}
\deleted{\uline{Exemple} Prenons $\varphi$ non dég.}
\[
\left|\begin{array}{l}
\varphi:\bigwedge^2V\longrightarrow k\qquad\text{non dégénérée}\\
\text{i.e. }\varphi\in\bigwedge^2\check V\\
\deleted{\unsure{\text{d'où }\check\varphi\in\bigwedge^2V}}\\
\chi=\varphi^2\in\bigwedge^4\check V\qquad\chi:\bigwedge^4V\longrightarrow k
\end{array}\right.
\]
\deleted{Alors $\chi$ comme forme sur $\bigwedge^2V$ est \unsure{non}}
\figure{\begin{tikzpicture}[x=0.3cm,y=0.3cm]\draw (0,4) .. controls (0.7,4) and (0.5,2) .. (1.4,1.2) .. controls (2.3,0.8) and (3.8,1.1) .. (4,0.1)--(4.3,-0.1);\end{tikzpicture}}

\uline{Ex} Prenons $\chi'$ non dégénérée, \deleted{\unsure{comme isom.}} $\alpha=\mathrm{id}$ i.e. $\widetilde W=\bigwedge^2V$.

Alors un choix convenable de $\xi\in\widetilde W=\bigwedge^2V$ donne n'importe quelle $\varphi$ voulue, \unsure{au facteur scalaire}, donc n'importe quelle $\check\varphi$ \inserted{$\in\bigwedge^2V$} non dégénérée, \unsure{au facteur scalaire}. Considérons alors $\check\varphi^2$ comme élément de \deleted{$\bigwedge^2V\unsure{\otimes}$} l'algèbre $\bigwedge^*V\otimes\bigwedge^*V$ [en écrivant $\bigwedge^2V\subset V\otimes V$] et $\check\varphi^2$ comme élément de $\bigwedge^2V\otimes\bigwedge^2V\subset\bigwedge^*V\otimes\bigwedge^*V$; le carré de $\check\varphi$ dans $\bigwedge^*V$ \inserted{(quand $\in\bigwedge^2V$)} est l'image de $\check\varphi^{(2)}$ précédent par $\bigwedge^2V\otimes\bigwedge^2V\longrightarrow\bigwedge^4V$. Utilisant la forme $\unsure{Q:Q_{\chi'}|Q=Q_{\chi'}}$ définie par $\chi'$, on voit que $\check\varphi^{(2)}$ définit un endom de $\bigwedge^2V$, \uline{soit $u$}. Je dis que l'orbite sous $u$ des éléments $\xi$ et $\check\varphi$ de $\bigwedge^2V$ \unsure{pourront être infinis|pourront être infinies}, si $\dim V\geq4$ (ou $6$?) ($V$ est de dim paire) et \deleted{\unsure{de}} \inserted{\unsure{linéairement engendrées et pleines}} de dim non bornée (\unsure{quand} $\dim V$ \unsure{croissant}\ldots).
\deleted{$\unsure{\xi}$}
\mspage{28}
